\documentclass[12pt]{article}
\PassOptionsToPackage{table}{xcolor}
\usepackage[utf8]{inputenc}
\usepackage[T1]{fontenc}
\usepackage[italian]{babel}
\usepackage{multicol}
\usepackage{enumitem}
\usepackage{array}
\usepackage{xcolor}
\usepackage{graphicx}
\usepackage{tikz}
\usetikzlibrary{
    shapes.geometric,
    arrows.meta,
    positioning,
    calc,
    patterns,
    decorations.pathmorphing
}
\usepackage[margin=2cm]{geometry}
\usepackage{amsmath,amssymb}
\setlength{\parindent}{0pt}
\definecolor{headercolor}{RGB}{220,220,220}
\newcounter{quesito}
\setcounter{quesito}{0}
\newcommand{\nuovoquesito}{\stepcounter{quesito}\textbf{Domanda \arabic{quesito}.}}
\newcommand{\itemdomanda}[2]{%
    \par\noindent
    \begin{minipage}[t]{\linewidth}
        \nuovoquesito \\ #1
        \begin{enumerate}[label=(\Alph*), nosep, leftmargin=*]
            #2
        \end{enumerate}
    \end{minipage}
    \vspace{10pt}
}
\begin{document}
\section*{Logica}
\hfill Data: 16 apr 2026

\itemdomanda{Basandosi sulle seguenti informazioni:
    \begin{enumerate}
        \item Gli unici animali di questa casa sono gatti.
        \item Ogni animale che ami guardare la luna è adatto ad essere un animale domestico.
        \item Quando detesto un animale, lo evito.
        \item Nessun animale che non si aggiri furtivamente nella notte è carnivoro.
        \item Non c'è gatto che non uccida i topi.
        \item Nessun animale si affeziona a me, fatta eccezione per quelli che sono in casa.
        \item I canguri non sono adatti ad essere animali domestici.
        \item Solo i carnivori uccidono i topi.
        \item Detesto gli animali che non si affezionano a me.
        \item Gli animali che si aggirano furtivamente nella notte amano sempre guardare la luna.
    \end{enumerate}
    Dunque:}{
    \item Evito i canguri
    \item I gatti non amano guardare la luna
    \item I canguri si affezionano a me
    \item Tutti gli animali carnivori sono in casa
    \item Non evito i gatti perché non sono carnivori
}

\itemdomanda{Date le seguenti premesse:
    \begin{enumerate}
        \item Tutti i colibrì hanno un piumaggio coloratissimo.
        \item Nessun uccello grosso si nutre di miele.
        \item Gli uccelli che non si nutrono di miele hanno colori smorti.
    \end{enumerate}
    Ne consegue che:}{
    \item I colibrì sono uccelli piccoli
    \item I colibrì sono uccelli grossi
    \item Gli uccelli con colori smorti sono colibrì
    \item Nessun uccello piccolo si nutre di miele
    \item Tutti gli uccelli grossi hanno il piumaggio colorato
}

\itemdomanda{Basandosi sulle seguenti regole:
    \begin{enumerate}
        \item Gli unici cibi che il mio medico mi concede sono quelli leggeri.
        \item Nessuno dei cibi che mi giovano è inadatto per la cena.
        \item Le torte nuziali sono sempre molto sostanziose.
        \item Il mio medico mi concede ogni cibo adatto per la cena.
    \end{enumerate}
    Quale affermazione è corretta?}{
    \item Le torte nuziali non mi giovano
    \item Le torte nuziali sono adatte per la cena
    \item Il mio medico mi concede le torte nuziali
    \item I cibi sostanziosi mi giovano sempre
    \item Tutti i cibi leggeri sono inadatti per la cena
}

\itemdomanda{Individuare il diagramma che rappresenta correttamente la relazione tra: \textbf{Animali}, \textbf{Mammiferi}, \textbf{Balene}.
\begin{center}
\begin{tikzpicture}[scale=0.25]
    % --- DIAGRAMMA 1 (Concentrici) ---
    \node[below] at (0,-2) {\textbf{1}};
    \draw (0,0) circle (2.2); % Animali
    \draw (0,0) circle (1.4); % Mammiferi
    \draw (0,0) circle (0.7); % Balene
    
    % --- DIAGRAMMA 2 (Intersecati) ---
    \begin{scope}[xshift=8cm]
    \node[below] at (0,-2) {\textbf{2}};
    \draw (0,0.8) circle (1.3); \draw (-0.8,-0.5) circle (1.3); \draw (0.8,-0.5) circle (1.3);
    \end{scope}

    % --- DIAGRAMMA 3 (Disgiunti in uno) ---
    \begin{scope}[xshift=16cm]
    \node[below] at (0,-2) {\textbf{3}};
    \draw (0,0) circle (1.8); \draw (-0.6,0) circle (0.8); \draw (0.6,0) circle (0.8);
    \end{scope}
\end{tikzpicture}
\end{center}}{
    \item DIAGRAMMA 2
    \item DIAGRAMMA 1
    \item DIAGRAMMA 3
    \item Nessuno dei precedenti
    \item Tutti i diagrammi sono corretti
}

\itemdomanda{Tre società: A, B e C, di cui almeno due sono italiane. Sapendo che se A è italiana lo è anche B, e che se C è italiana lo è anche A, e che tra B e C almeno una non è italiana, possiamo dedurre che:}{
    \item C non è italiana e B è italiana
    \item A non è italiana e B è italiana
    \item A, B e C sono italiane
    \item A e C sono italiane
    \item C è italiana e B non è italiana
}

\itemdomanda{Tre amiche, Maria, Carla e Daniela, vanno spesso a teatro. Si sa che se Daniela va a teatro, ci va anche Maria; Maria e Carla non ci vanno mai tutte e due insieme; può anche succedere che Daniela o Carla vadano a teatro da sole; però, se Carla va a teatro, allora ci va anche Daniela. Supponendo vere tali premesse, chi è andata al cinema in coppia?}{
    \item Maria e Carla
    \item Nessuna delle tre
    \item Maria e Daniela
    \item Carla e Daniela
    \item Ciascuna di loro
}

\end{document}
